% Sample output for slide-script-tex-generator
% Style: top-slide-manuscript (default)
% Engine: XeLaTeX
% Put slides.pdf and this .tex file in the same folder.

\documentclass[11pt,a4paper]{article}

\usepackage[margin=1.5cm]{geometry}
\usepackage{graphicx}
\usepackage{fontspec}
\usepackage{xcolor}
\usepackage{parskip}

\newcommand{\SlidePDF}{slides.pdf}
\newcommand{\DocTitle}{Presentation Manuscript}

\newenvironment{SlideManuscriptBlock}[3][normalsize]{
  \section*{#2}
  \begin{center}
    \includegraphics[page=#3,width=0.92\textwidth]{\SlidePDF}
  \end{center}
  \begingroup
  \csname #1\endcsname
}{
  \endgroup
  \newpage
}

\begin{document}
\section*{\DocTitle}

\begin{SlideManuscriptBlock}[large]{Slide 1}{1}
Opening remarks and agenda overview. This page has very short script text, so a larger manuscript size is used.
\end{SlideManuscriptBlock}

\begin{SlideManuscriptBlock}[normalsize]{Slide 2}{2}
This slide contains the core background context, assumptions, and transition sentences for the next section. The content length is moderate, so normalsize is used.
\end{SlideManuscriptBlock}

\begin{SlideManuscriptBlock}[small]{Slide 3}{3}
This section has a longer manuscript body with multiple key points, supporting evidence, and closing notes. A smaller size helps prevent overflow while keeping the full script on the page in a stable and readable way.
\end{SlideManuscriptBlock}

\end{document}
